\documentclass[11pt,a4paper]{article}
\usepackage[utf8]{inputenc}
\usepackage[T1]{fontenc}
\usepackage{lmodern}
\usepackage[a4paper,margin=1in]{geometry}
\usepackage{amsmath,amssymb}
\usepackage{enumitem}
\usepackage{hyperref}
\setlist[itemize]{leftmargin=1.5em,itemsep=0.25em,topsep=0.25em}

\title{\textbf{Initial Reading Notes on SABR Models for Synthetic FX Rates}}
\author{Master 2 El Karoui --- 2-day progress note}
\date{\today}

\begin{document}
\maketitle
\vspace{-1.0em}

\section*{1. Objective of the paper}
\begin{itemize}
    \item The paper studies how to infer cross-FX option smiles from two liquid FX smiles quoted versus a common base currency.
    \item Main question: if \(\widetilde X_{12}=\widetilde X_{10}/\widetilde X_{20}\), can we obtain a reliable SABR smile for \(\widetilde X_{12}\) from the SABR models of \(\widetilde X_{10}\) and \(\widetilde X_{20}\)?
    \item Main result: yes, through an effective lognormal SABR approximation accurate to \(O(\varepsilon^2)\).
\end{itemize}

\section*{2. Global understanding}
\begin{itemize}
    \item \textbf{Main text:} financial setup, cross construction, and practical formulas for effective SABR parameters.
    \item \textbf{Appendix A:} general asymptotic reduction from a multi-factor stochastic-vol model to an effective one-factor forward equation.
    \item \textbf{Appendix B:} specialization of Appendix A to cross-FX (\(n=2\), signed loadings), giving explicit cross formulas.
    \item \textbf{Overall strategy:} derive marginal density PDE in one dimension, then match it to SABR's effective forward PDE.
\end{itemize}

\section*{3. Part I focused on}
\begin{itemize}
    \item I chose to focus mainly on \textbf{Appendix A}.
    \item Reason: it is the mathematical core that explains \emph{why} an effective SABR emerges, not only \emph{what} formulas to use.
    \item This section gives the reduction pipeline: high-dimensional dynamics \(\to\) closed marginal PDE \(\to\) effective SABR parameters.
\end{itemize}

\section*{4. What I understood well}
\begin{itemize}
    \item Cross rate dynamics comes naturally from It\^o on \(\widetilde X_{12}=\widetilde X_{10}/\widetilde X_{20}\), with drift removed in the appropriate forward measure.
    \item The cross is driven by a signed combination of component drivers; correlation \(c\) controls cancellation/amplification of noise.
    \item In Appendix A, the key closure issue is clear:
    \[
    \partial_T Q^{(0)}=\frac12\varepsilon^2\partial_{XX}\!\bigl(X^2Q^{(2)}\bigr),
    \]
    so one needs a constitutive relation \(Q^{(2)}\) in terms of \(Q^{(0)}\).
    \item The effective-SABR parameters are obtained by matching the resulting one-dimensional PDE with SABR's asymptotic forward PDE.
\end{itemize}

\section*{5. Key intuitions}
\begin{itemize}
    \item \textbf{Cross FX dynamics:} ratio structure implies a difference of component return shocks.
    \item \textbf{Effective volatility:} instantaneous variance has a quadratic form combining both vols and their correlation.
    \item \textbf{Why effective SABR emerges:} after asymptotic reduction, the marginal PDE has the same corrected structure as SABR (\(1+2\varepsilon \cdot +\varepsilon^2\cdot\)).
    \item \textbf{Purpose of asymptotic reduction:} preserve pricing-relevant information while collapsing dimensionality.
\end{itemize}

\section*{6. What I am still working on}
\begin{itemize}
    \item Some technical details in the variable transforms used to prove the \(Q^{(2)}\)-\(Q^{(0)}\) relation.
    \item Detailed interpretation of all higher-order coefficients (especially mixed-correlation contributions).
    \item Full comparison between general-correlation and single-stochastic-vol cases in Appendix B.
    \item Calibration implications across maturities, since the mapped \(\alpha_{\mathrm{eff}}\) is expiry-dependent.
\end{itemize}

\end{document}
